%% LyX 2.5.1 created this file.  For more info, see https://www.lyx.org/.
%% Do not edit unless you really know what you are doing.
\documentclass[english,hebrew]{article}
\usepackage[HE8,T1]{fontenc}
\usepackage[utf8]{inputenc}
\usepackage{babel}
\usepackage{amstext}
\usepackage{amssymb}
\usepackage[unicode=true]
 {hyperref}
\begin{document}
\title{מבוא לטופולוגיה קבוצתית: סגור, נקודות הצטברות ופונקציות רציפות}
\maketitle
\begin{description}
\item [{קטגוריות:}] טופולוגיה
\item [{תגים:}] מרחבים טופולוגיים, פונקציות רציפות
\item [{מזהה:}] \L{topological\_spaces\_closure\_and\_continuous\_functions}
\end{description}
לפני שנגיע לדבר העיקרי שאני רוצה לדבר עליו, \textbf{פונקציות רציפות},
בואו נפתח עם חימום שמשלים חוב שיש לי.

בפוסט הקודם ראינו את ההגדרה הבסיסית של מרחב טופולוגי. בואו ניזכר בה:

\textbf{מרחב טופולוגי} הוא זוג \L{$\left(X,\mathcal{T}\right)$} של
קבוצה \L{$X$} ו-\L{$\mathcal{T}\subseteq\mathcal{P}\left(X\right)$}
הוא אוסף של תת-קבוצות של \L{$X$} שאנחנו קוראים להן \textbf{קבוצות
פתוחות}. הדרישה מ-\L{$\mathcal{T}$} היא שיתקיימו שלוש התכונות הבאות:
\begin{itemize}
\item \L{$\emptyset,X\in\mathcal{T}$}
\item אם יש לנו אוסף קבוצות \L{$A_{i}\in\mathcal{T}$} אז \L{$\bigcup_{i}A_{i}\in\mathcal{T}$}
\item אם יש לנו מספר סופי של קבוצות \L{$A_{1},\ldots,A_{n}\in\mathcal{T}$}
אז \L{$\bigcap^{n}_{i=1}A_{i}\in\mathcal{T}$}
\end{itemize}
עכשיו, בפוסטים הקודמים ראינו כמה מושגים קשורים, למשל \textbf{קבוצות
סגורות}. אחת מהתוצאות שראינו עבור מרחבים מטריים הייתה שהמשלימה של
קבוצה סגורה היא פתוחה, וההפך. לכן מתבקש להפוך את זה להגדרה במקרה שלנו:
\textbf{קבוצה סגורה} היא קבוצה שהמשלימה שלה פתוחה. אבל איך זה מתקשר
לאופן שבו הגדרנו פעם קבוצה סגורה? בואו נזמן לרגע את רשימת ה\textquotedblright מה
עשינו פה\textquotedblright{} מהפוסט של הקבוצות הסגורות:
\begin{itemize}
\item בהינתן \L{$X$} כלשהי, לקחנו \L{$\mathcal{T}\subseteq\mathcal{P}\left(X\right)$}
וקראנו לאיברים שלה \textbf{קבוצות פתוחות}.
\item לכל \L{$x\in X$}, קראנו בשם \textbf{סביבה} של \L{$x$} לקבוצה פתוחה
\L{$U\in\mathcal{T}$} כלשהי כך ש-\L{$x\in U$}.
\item הגדרנו \textbf{נקודת הצטברות} של קבוצה \L{$A$} בתור נקודה \L{$x$}
שכל סביבה שלה כוללת איבר של \L{$A$} ששונה מ-\L{$x$}.
\item הגדרנו \textbf{קבוצה סגורה} בתור קבוצה שכוללת את כל נקודות ההצטברות
שלה.
\end{itemize}
אחר כך הראינו \textbf{במרחבים מטריים} את השקילות של ההגדרה הזו של
קבוצה סגורה וההגדרה של \textquotedblleft המשלימה שלה פתוחה\textquotedblright .
אבל אנחנו יכולים לעשות את זה באופן כללי, עם האקסיומות של מרחבים טופולוגיים.
וכדי לעשות את זה בצורה נחמדה שתלמד אותנו משהו חדש, אפשר להכניס לתמונה
מושגים חדשים שההגדרה שלהם בשפה טופולוגית היא טבעית למדי - \textbf{הסגור}
)\L{Closure}( וה\textbf{פנים} )\L{Interior}( של קבוצה.
\begin{itemize}
\item הסגור של קבוצה \L{$A$} שמסומן \L{$\overline{A}$} הוא חיתוך כל הקבוצות
הסגורות שמכילות את \L{$A$} )מכיוון ש-\L{$X$} היא סגורה כי המשלימה
שלה היא \L{$\emptyset$} הפתוחה, אז החיתוך הזה לא נלקח על אוסף ריק
של קבוצות, שזה כידוע דבר שאי אפשר לעשות(. כלומר, \L{$\overline{A}$}
היא הקבוצה הסגורה \textbf{המינימלית} שמכילה את \L{$A$}.
\item הפנים של קבוצה \L{$A$} שמסומן \L{$\text{int}A$} הוא איחוד כל הקבוצות
הפתוחות שמוכלות ב-\L{$A$}. כלומר, \L{$\text{int}A$} היא הקבוצה הפתוחה
\textbf{המקסימלית} שמוכלת ב-\L{$A$}.
\end{itemize}
מה זה אומר, בעצם? על פי הגדרה, \L{$\text{int}A\subseteq A\subseteq\overline{A}$}
)כי כשמאחדים קבוצות שכולן מוכלות ב-\L{$A$} גם התוצאה מוכלת ב-\L{$A$},
וכנ\textquotedblright ל כשחותכים קבוצות שכולן מכילות את כל \L{$A$}(.
אז הרעיון הוא שכשעוברים מ-\L{$A$} אל \L{$\overline{A}$} מוסיפים
לה נקודות וכשעוברים אל \L{$\text{int}A$} מורידים ממנה נקודות, אבל
איזה נקודות? אינטואיטיבית, נקודות השפה של \L{$A$}, מה שקראנו לו \L{$\partial A$}
- כל הנקודות שבכל סביבה שלהן יש איבר מ-\L{$A$} ואיבר מ-\L{$A^{c}$},
המשלימה של \L{$A$}.

בואו נראה את זה פורמלית. ראשית אני אוכיח ש-\L{$\overline{A}=A\cup\partial A$}.
בואו נתחיל עם נקודה \L{$x\in\partial A$} ונוכיח שהיא ב-\L{$\overline{A}$}.
בשביל זה מספיק להראות שלכל קבוצה סגורה \L{$D$} כך ש-\L{$A\subseteq D$},
\L{$x\in D$}. נניח בשלילה שזה לא המצב, כלומר \L{$x\notin D$}, כלומר
\L{$x\in D^{c}\subseteq A^{c}$}. מכיוון ש-\L{$D$} סגורה, המשלימה
שלה \L{$D^{c}$} פתוחה ולכן \L{$D^{c}$} היא סביבה של \L{$x$} שמוכלת
כולה ב-\L{$A^{c}$} ולכן לא מכילה אף נקודה מ-\L{$A$}, בסתירה לכך
ש-\L{$x\in\partial A$}, אז הנחת השלילה ש-\L{$x\notin D$} הייתה שגויה
וקיבלנו \L{$x\in D$}, מה שמסיים את הכיוון הזה של ההוכחה.

עכשיו בואו ניקח נקודה \L{$x\in\overline{A}$}, נניח שהיא לא ב-\L{$A$}
)כי אם היא כן, סיימנו( ונוכיח שהיא ב-\L{$\partial A$}. ניקח סביבה
\L{$U$} של \L{$x$}. מכיוון ש-\L{$x$} לא ב-\L{$A$}, הסביבה הזו
אוטומטית כוללת איבר מ-\L{$A^{c}$}, נשאר רק להראות שיש \L{$a\in A$}
ששייך ל-\L{$U$}. אז בואו נניח בשלילה שאין כזה, כלומר \L{$U\subseteq A^{c}$},
ולכן \L{$A\subseteq U^{c}$}. עכשיו כבר אפשר לראות לאן זה הולך: סביבה
היא בהגדרה קבוצה פתוחה, ולכן \L{$U^{c}$} היא קבוצה סגורה, והיא אחת
מהקבוצות הסגורות שמכילות את \L{$A$} ולכן מאלו שהחיתוך שלהן נותן את
\L{$\overline{A}$}. אז קיבלנו \L{$x\in\overline{A}\subseteq U^{c}$}
בסתירה לכך ש-\L{$x\in U$}, וסיימנו את ההוכחה.

ההוכחה הזו ממחישה יפה את הקושי שיש לפעמים בטופולוגיה. כל צעד כאן הוא
פשוט ודי מובן בפני עצמו, אבל לי אישית, למרות שזה עתה כתבתי את ההוכחה
הזו מהראש, אין אינטואיציה טובה מה הולך כאן. הכל מרגיש לי כמו משחקים
בהגדרות. לא שיש משהו רע בזה... אבל בלי העוגן של מרחבים מטריים, אפשר
להבין כמה קל ללכת לאיבוד לפעמים כשהעניינים מסתבכים.

אוקיי, מה עם \L{$\text{int}A$}? אנחנו רוצים להוכיח ש-\L{$\text{int}A=A\backslash\partial A$},
כלומר שהפנים של \L{$A$} כולל רק את אברי \L{$A$} שאינם ב-\L{$\partial A$}.
דרך אחרת לנסח את זה היא בתור \L{$\text{int}A=A\cap\left(\partial A\right)^{c}$},
כלומר הפנים רק את אברי \L{$A$} שבנוסף לכך גם לא שייכים ל-\L{$\partial A$}.

הציפייה היא שההוכחה תהיה איכשהו \textquotedblleft דואלית\textquotedblright{}
למה שכבר ראינו, הרי כל ההגדרות הן מעין גרסת מראה אחת של השניה, אבל
אם ככה - אולי נמצא דרך קיצור ונשתמש במה שכבר הוכחנו, ואולי נלמד דברים
חדשים על הדרך?

הטריק הוא באמת ללכת אל הדואלי - כלומר, לקחת משלימים לקבוצות שעובדים
איתן. אם \L{$\text{int}A$} הוא איחוד של קבוצות פתוחות שמוכלות ב-\L{$A$},
כלומר \L{$\text{int}A=\bigcup_{U\subseteq A}U$}, אז אם ניקח משלים
לשני האגפים ונשתמש בכללי דה-מורגן, נקבל

\L{$\left(\text{int}A\right)^{c}=\bigcap_{U\subseteq A}U^{c}$}

עכשיו, אם \L{$U\subseteq A$} אז \L{$A^{c}\subseteq U^{c}$} )נסו
להוכיח את זה! זה קל!( כך שאפשר לכתוב את זה גם בתור

\L{$\left(\text{int}A\right)^{c}=\bigcap_{A^{c}\subseteq U^{c}}U^{c}$}

ובגלל שהמשלים של קבוצה פתוחה הוא קבוצה סגורה, אנחנו מקבלים ש-\L{$\left(\text{int}A\right)^{c}=\overline{A^{c}}$}-
המשלים של הפנים של \L{$A$} הוא הסגור של המשלים של \L{$A$} )דואליות
בפעולה!!!!! אגב, אני לא סתם משתמש במילה הזו, בתורת הקטגוריות יש מהומה
שלמה סביב הקונספט הזה(.

עכשיו אפשר להשתמש במה שאנחנו כבר יודעים: \L{$\left(\text{int}A\right)^{c}=\overline{A^{c}}=A^{c}\cup\partial A^{c}$}.
לכן, אם ניקח משלים לשני האגפים ונשתמש שוב בדה-מורגן, נקבל

\L{$\text{int}A=A\cap\left(\partial A^{c}\right)^{c}$}

עכשיו, מה זה \L{$\partial A^{c}$}? זה כל הנקודות שבכל סביבה שלהן
יש נקודה מ-\L{$A^{c}$} ונקודה מ-\L{$\left(A^{c}\right)^{c}=A$}.
זו... אותה הגדרה בדיוק כמו של \L{$\partial A$}, אז באופן לא כל כך
מפתיע, השפה של \L{$A$} והשפה של \L{$A^{c}$} הן אותו דבר, \L{$\partial A=\partial A^{c}$},
אז קיבלנו

\L{$\text{int}A=A\cap\partial A^{c}$}

וזה בדיוק מה שרצינו. יפה!

עכשיו אפשר להסיק כמה מסקנות. ראשית, \L{$A$} סגורה אם ורק אם \L{$\overline{A}=A$}.
למה? בכיוון אחד, \L{$\overline{A}$} סגורה כי היא חיתוך של קבוצות
סגורות, אז אם \L{$A=\overline{A}$}, בוודאי ש-\L{$A$} סגורה. בכיוון
השני, אם \L{$A$} סגורה אז היא עצמה חלק מהחיתוך שנותן את \L{$\overline{A}$}
וכל קבוצה בחיתוך מכילה את \L{$A$} אז כל מה שנשאר מהחיתוך הוא \L{$A$}.
באותו אופן מראים ש-\L{$A$} פתוחה אם \L{$A=\text{int}A$}.

עכשיו בואו נחזור לנקודות הצטברות. הנה שוב ההגדרה, כדי שתהיה מול העיניים:
\begin{itemize}
\item \textbf{נקודת הצטברות} של קבוצה \L{$A$} היא נקודה \L{$x$} שכל סביבה
שלה כוללת איבר של \L{$A$} ששונה מ-\L{$x$}.
\end{itemize}
אנחנו רוצים להראות ש-\L{$A$} סגורה אם ורק אם היא כוללת כל נקודת הצטברות
שלה. אם \L{$A$} סגורה, אז \L{$A=\overline{A}$} ולכן \L{$A$} כוללת
את כל \textbf{השפה} שלה - ונקודת הצטברות של \L{$A$} או שייכת אל \L{$A$}
או שהיא בשפה של \L{$A$}, כי אם \L{$x\notin A$} היא נקודת הצטברות
אז כל סביבה שלה כוללת איבר של \L{$A$} )כי \L{$x$} נקודת הצטברות(
וגם איבר שאינו של \L{$A$} )את \L{$x$} עצמו, כי \L{$x\notin A$}(.
זה נותן כיוון אחד. בכיוון השני, נניח ש-\L{$A$} כוללת את כל נקודות
ההצטברות שלה, אז היא בוודאי תכלול גם את השפה שלה כי כל נקודת שפה שאינה
ב-\L{$A$} היא גם נקודת הצטברות, כי כל סביבה שלה כוללת נקודה ב-\L{$A$}
ואם היא עצמה לא ב-\L{$A$} הנקודה הזו שונה ממנה. זה מסיים את ההוכחה
הזו.

עכשיו אפשר להשלים את רשימת המושגים שהגדרנו בהקשר של מרחבים מטריים
ולהגדיר אותם על טהרת המרחבים הטופולוגיים. מה שהתחלתי ממנו היה \textbf{גבולות}
ובשלב הזה אנחנו כבר יודעים איך לנסח אותם בלי מטריקות:
\begin{itemize}
\item אם \L{$X$} הוא מרחב טופולוגי ו-\L{$a_{1},a_{2},a_{3},\ldots\in X$}
היא סדרת איברים של \L{$X$}, אומרים שהסדרה מתכנסת לגבול \L{$a$} ומסמנים
זאת \L{$\lim_{n\to\infty}a_{n}=a$} אם לכל סביבה \L{$U$} של \L{$a$}
קיים \L{$N$} כך שלכל \L{$n>N$}, \L{$a_{n}\in U$}
\item אם \L{$X,Y$} הם מרחבים טופולוגיים ו-\L{$f:X\to Y$} פונקציה, אומרים
ש-\L{$f\left(x\right)$} שואפת לגבול \L{$L\in Y$} כאשר \L{$x$} שואף
אל \L{$x_{0}\in X$} ומסמנים זאת \L{$\lim_{x\to x_{0}}f\left(x\right)=L$}
אם לכל סביבה \L{$V\subseteq Y$} של \L{$L$} קיימת סביבה \L{$U\subseteq X$}
של \L{$x_{0}$} כך ש-\L{$f\left(x\right)\in V$} לכל \L{$x\in U$}
כך ש-\L{$x\ne x_{0}$}.
\end{itemize}
שני אלו הם תרגום כמעט ישיר של הגדרת הגבול ה\textquotedblright רגילה\textquotedblright .
בגבול רגיל אנחנו לא מדברים על סביבות של \L{$x$} אלא על כל הנקודות
במרחק כלשהו מ-\L{$x$}, וכבר ראינו ש\textquotedblright סביבה\textquotedblright{}
היא ההכללה הטבעית של זה.

עכשיו אפשר להגיע אל המושג המרכזי שחתרנו אליו - \textbf{רציפות}. ההגדרה
הרגילה של רציפות היא פשוט \L{$\lim_{x\to x_{0}}f\left(x\right)=f\left(x_{0}\right)$}
שזה מה שמתקבל כשלוקחים את הגדרת הגבול הרגילה של פונקציה ומורידים את
הדרישה \L{$x\ne x_{0}$} )ולוקחים \L{$L=f\left(x_{0}\right)$}(.
\begin{itemize}
\item אם \L{$X,Y$} הם מרחבים טופולוגיים ו-\L{$f:X\to Y$} פונקציה, אומרים
ש-\L{$f\left(x\right)$} \textbf{רציפה} בנקודה \L{$x_{0}$} אם לכל
סביבה \L{$U\subseteq Y$} של \L{$f\left(x_{0}\right)$} קיימת סביבה
\L{$V\subseteq X$} של \L{$x_{0}$} כך ש-\L{$f\left(x\right)\in U$}
לכל \L{$x\in V$}.
\end{itemize}
אבל אפשר לעשות משהו מעניין יותר. כפי שהיא מוגדרת כאן, רציפות היא תכונה
\textbf{נקודתית}. אינטואיטיבית, היא אומרת שהפונקציה \textquotedblleft מבטיחה
וגם מקיימת\textquotedblright{} משהו בנקודה \L{$x_{0}$}: הגבול יוצר
את הרושם, ככל שמתקרבים אל \L{$x_{0}$}, שהפונקציה הולכת להחזיר ערך
ספציפי בנקודה \L{$x_{0}$} - והפלא ופלא, כשמגיעים לשם מתברר שהערך
הזה הוא אכן \L{$f\left(x_{0}\right)$}. אבל אפשר גם לדבר על רציפות
בתור תכונה של פונקציה בתחום שלם מסוים, או אפילו על כל המרחב. אפשר
לנסות לומר ש-\L{$f:X\to Y$} היא \textbf{רציפה} אם היא רציפה ב-\L{$x_{0}$}
לכל \L{$x_{0}\in X$}, אבל אפשר לנסח את זה גם אחרת.

ראשית, בואו נסתכל על הניסוח של רציפות. אני אומר בסוף \textquotedblleft כך
ש-\L{$f\left(x\right)\in V$} לכל \L{$x\in U$}\textquotedblright .
יש דרך אחרת לנסח את הקריטריון הזה בעזרת מושג \textbf{המקור} של קבוצה.
אם \L{$f:X\to Y$} היא פונקציה ו-\L{$V\subseteq Y$}, אז \textbf{המקור}
של \L{$V$} תחת \L{$f$} הוא הקבוצה
\begin{itemize}
\item \L{$f^{-1}\left(V\right)\triangleq\left\{ x\in X\ |\ f\left(x\right)\in V\right\} $}
\end{itemize}
כלומר, כל האיברים ב-\L{$X$} שתחת הפונקציה \L{$f$} מחזירים משהו ב-\L{$V$}.
שימו לב שה-\L{$f^{-1}$} שמופיע פה הוא \textbf{סימון}; בדרך כלל משתמשים
ב-\L{$f^{-1}$} כדי לתאר את הפונקציה ההופכית של \L{$f$}, אבל פונקציה
הופכית כזו קיימת רק אם \L{$f$} חח\textquotedblright ע ועל, ואני לא
דורש את זה פה. אז הקונבנציה היא שכשכותבים \L{$f^{-1}\left(V\right)$}
הכוונה למקור - וקשה להתבלבל כי מבחינה פורמלית אפילו אם \L{$f^{-1}$}
הייתה קיימת, אמורים להפעיל אותה על איברים ולא על קבוצות אם כי כמובן
שגם \L{$f\left(V\right)\triangleq\left\{ f\left(x\right)\ |\ x\in V\right\} $}
הוא סימון מקובל עבור \textbf{התמונה} של \L{$V$} תחת \L{$f$}.

עכשיו, להגיד \textquotedblleft כך ש-\L{$f\left(x\right)\in V$} לכל
\L{$x\in U$}\textquotedblright{} זה בעצם לומר \L{$x\in f^{-1}\left(V\right)$}
לכל \L{$x\in U$}, כלומר \L{$U\subseteq f^{-1}\left(V\right)$}. לכן
בניסוח בעזרת מקור של ההגדרה של רציפות, \L{$f\left(x\right)$} רציפה
בנקודה \L{$x_{0}$} אם לכל סביבה \L{$V\subseteq Y$} של \L{$f\left(x_{0}\right)$}
קיימת סביבה \L{$U\subseteq X$} של \L{$x_{0}$} כך ש-\L{$U\subseteq f^{-1}\left(V\right)$}.
האם אפשר לפשט את זה עוד? ובכן, כן.

ראשית, בואו נשים לב לכך שאם היה מובטח לי ש-\L{$f^{-1}\left(V\right)$}
היא קבוצה \textbf{פתוחה}, אז הקריטריון היה מתקיים מאליו: \L{$f^{-1}\left(V\right)$}
היא גם קבוצה פתוחה וגם מכילה את \L{$x_{0}$} )כי \L{$f\left(x_{0}\right)\in V$}
כי כל הרעיון הוא ש-\L{$V$} היא סביבה של \L{$f\left(x_{0}\right)$}(
ולכן אפשר לבחור \L{$U=f^{-1}\left(V\right)$} ולסיים בזה. אבל האם
\textbf{מובטח} לנו ש-\L{$f^{-1}\left(V\right)$} היא קבוצה פתוחה,
אם \L{$f$} רציפה?

ובכן, כן.

ברגע שאנחנו עוברים מרציפות נקודתית לרציפות בכל נקודה במרחב, אנחנו
כבר לא צריכים לחשוב על \L{$V$} בתור \textquotedblleft הסביבה של \L{$x_{0}$}\textquotedblright .
הקבוצה הפתוחה \L{$V$} היא הסביבה של \textbf{כל} הנקודות שבתוכה. ואפשר
להשתמש ברציפות עבור \textbf{כל} הנקודות בתוכה שיש להן מקור. לכל \L{$x\in f^{-1}\left(V\right)$}
קיימת, על פי הגדרת הרציפות, סביבה \L{$U_{x}$} כך ש-\L{$U_{x}\subseteq f^{-1}\left(V\right)$}.
לכן

\L{$\bigcup_{x\in f^{-1}\left(V\right)}U_{x}\subseteq f^{-1}\left(V\right)$}

)כי זה איחוד של קבוצות שכולן מוכלות ב-\L{$f^{-1}\left(V\right)$}(,
ומצד שני לכל \L{$x\in f^{-1}\left(V\right)$} קיימת קבוהצ \L{$U_{x}$}
באיחוד שמכילה את \L{$x$}, כך ש-\L{$f^{-1}\left(V\right)=\bigcup_{x\in f^{-1}\left(V\right)}U_{x}$}.
בפרט, קיבלנו ש-\L{$f^{-1}\left(V\right)$} הוא איחוד של סביבות, כלומר
קבוצות פתוחות, \L{$f^{-1}\left(V\right)$} היא קבוצה פתוחה - וזה היה
נכון עבור \textbf{כל} קבוצה פתוחה \L{$V$}, כי כל קבוצה פתוחה היא
הסביבה של כל הנקודות שלה )אם \L{$V$} לא הייתה פתוחה, היא לא הייתה
סביבה של אף נקודה(. זה מוביל אותנו להגדרה האלטרנטיבית הבאה לרציפות,
שהיא כנראה ההגדרה המרכזית בטופולוגיה:
\begin{itemize}
\item \L{$f:X\to Y$} היא \textbf{רציפה} אם לכל קבוצה פתוחה \L{$V\subseteq Y$},
המקור \L{$f^{-1}\left(V\right)$} הוא קבוצה פתוחה
\end{itemize}
ההגדרה הזו עשויה להרגיש קצת מוזרה בהתחלה. על פניו מרגיש טבעי יותר
להגדיר ההפך, לא? שאם \L{$U\subseteq X$} היא קבוצה פתוחה, אז \L{$f\left(U\right)$}
גם תהיה פתוחה? אבל לא, לכן לדעתי כדאי לצאת מההגדרה ה\textquotedblright רגילה\textquotedblright{}
לרציפות - אנחנו רואים איך בהגדרה הזו \L{$f^{-1}$} באה מאליה. אבל
זה לא אומר שהתכונה השניה לא חשובה בפני עצמה, ובפרט היא חוזרת לגמרי
למשחק כשאנחנו מגדירים את מה שנקרא \textbf{הומיאומורפיזם} שנראה לי
שכדאי להציג כבר עכשיו.

כשמתחילים ללמוד מתמטיקה, מהר מאוד נתקלים בכל מני פונקציות שנקראות
\textbf{מורפיזמים}. הרעיון הכללי שמאחורי השם הזה הוא פונקציה בין שני
מרחבים שיש להם \textbf{מבנה} כלשהו והפונקציה \textbf{משמרת} את המבנה
הזה )ואז מגיעה \textbf{תורת הקטגוריות} ומכלילה את כל הסיפור אבל לא
ניכנס לזה הפעם(. בפרט נתקלים בשם \textbf{איזומורפיזם} כדי לתאר פונקציה
שהיא חד-חד-ערכית, על, ומשמרת את המבנה. עכשיו, הקטע בפונקציות חד-חד-ערכיות
ועל בין שתי קבוצות \L{$A,B$} הוא שהן מאפשרות לנו לחשוב על \L{$B$}
בתור \textquotedblleft\L{$A$} אחרי ששינינו שמות לאיברים שלה\textquotedblright{}
)יש לי פוסט על זה \href{https://gadial.net/2020/05/15/invertible_functions/}{כאן}(.
זה כבר לא עובד עבור פונקציות חח\textquotedblright ע ועל כלליות אם
יש על \L{$A,B$} מבנה כלשהו - למשל אם הם \textbf{מרחבים וקטוריים}
או \textbf{קבוצות סדורות} או \textbf{גרפים} או \textbf{חבורות} )כל
אלו הן סיטואציות שבהן אפשר להיתקל באיזומורפיזמים כבר בתחילת לימודי
המתמטיקה(. אם מוסיפים את הדרישה הנוספת, של שימור המבנה, אז מקבלים
שוב את נקודת המבט של \textquotedblleft\L{$A$} ו-\L{$B$} הם אותו
דבר עד כדי שינוי שמות\textquotedblright . בטופולוגיה המילה שבאה לתאר
פונקציות כאלו שהן חח\textquotedblright ע, על ומשמרות מבנה היא \textbf{הומיאומורפיזם}.

מה ה\textquotedblright מבנה\textquotedblright{} שרוצים לשמר הפעם? אמרנו
שמרחב טופולוגי הוא זוג \L{$\left(X,\mathcal{T}_{X}\right)$} כך ש-\L{$X$}
היא בסך הכל קבוצה, ואילו \L{$\mathcal{T}$} הוא אוסף של תת-קבוצות
של \L{$X$} שנקראות \textbf{קבוצות פתוחות}. אם \L{$f$} היא הומיאומורפיזם,
אז בנוסף לכך שהיא חח\textquotedblright ע ועל היינו רוצים שהיא תשמר
קבוצות פתוחות, כלומר אם \L{$U\subseteq X$} היא פתוחה, שגם \L{$f\left(U\right)\subseteq Y$}
תהיה פתוחה )פורמלית, אם \L{$U\in\mathcal{T}_{X}$} אז \L{$f\left(U\right)\in\mathcal{T}_{Y}$}(,
אבל זה לבדו לא מספיק, צריך עוד כיוון: שאם \L{$U$} \textbf{איננה}
פתוחה גם \L{$f\left(U\right)$} לא תהיה כזו )ולכן אינטואיטיבית אפשר
לחשוב על זה כאילו מרחיבים באופן טבעי את \L{$f$} לפונקציה חח\textquotedblright ע
ועל \L{$f:\mathcal{T}_{X}\to\mathcal{T}_{Y}$}(. עכשיו, אם \L{$f$}
היא חח\textquotedblright ע ועל ו-\L{$V\subseteq Y$} היא קבוצה כלשהי,
אז \L{$f^{-1}\left(V\right)$} זו אותה קבוצה בדיוק בין אם חושבים עליה
בתור \textbf{המקור} של \L{$V$} תחת \L{$f$} במובן שהראיתי קודם, או
בתור \textbf{התמונה} של \L{$f^{-1}$} )שעכשיו היא פונקציה אמיתית ולא
סתם סימון(, כלומר

\L{$f^{-1}\left(V\right)=\left\{ x\in X\ |\ f\left(x\right)\in V\right\} =\left\{ f^{-1}\left(y\right)\ |\ y\in V\right\} $}

זו טענה שצריך להוכיח אבל היא פשוטה )תרגיל!(. עבורנו היא מועילה כי
עכשיו קל לנסח את התכונה הנוספת שאנחנו מחפשים. אמרתי שאני רוצה לדרוש
את \textquotedblleft אם \L{$U$} איננה פתוחה גם \L{$f\left(U\right)$}
איננה פתוחה\textquotedblleft . הטענה הזו שקולה לוגית אל \textquotedblleft אם
\L{$f\left(U\right)$} פתוחה גם \L{$U$} פתוחה\textquotedblright .
עכשיו, אם נסמן \L{$V=f\left(U\right)$}, אז \L{$f^{-1}\left(V\right)=f^{-1}\left(f\left(U\right)\right)=U$}
)זו טענה שצריך להוכיח אבל היא פשוטה(, כלומר אפשר לנסח את הדרישה שלנו
בתור \textquotedblleft אם \L{$V\subseteq Y$} פתוחה גם \L{$f^{-1}\left(V\right)$}
פתוחה\textquotedblright . וזו ההגדרה לרציפות שכבר ראינו. אז אפשר להגדיר
ש-
\begin{itemize}
\item \L{$f:X\to Y$} היא הומואימורפיזם אם היא חח\textquotedblright ע, על,
רציפה ומעבירה קבוצות פתוחות לקבוצות פתוחות
\end{itemize}
אפשר לנסח את זה אפילו עוד יותר אלגנטי. במקום לומר \textquotedblleft חח\textquotedblright ע
ועל\textquotedblright{} מספיק לומר \textbf{הפיכה}. ואם לכל \L{$U$}
פתוחה גם \L{$f\left(U\right)$} אז אפשר לנסח את זה גם בתור \textquotedblleft לכל
\L{$U$} פתוחה גם \L{$\left(f^{-1}\right)^{-1}\left(U\right)$} פתוחה\textquotedblright ,
כלומר- \textbf{המקור} של \L{$U$} על פי הפונקציה \L{$f^{-1}$} היא
קבוצה פתוחה - מה שאומר ש-\L{$f^{-1}$} היא \textbf{רציפה}. אז אפשר
לומר
\begin{itemize}
\item \L{$f:X\to Y$} היא הומואימורפיזם אם היא הפיכה, רציפה, וההופכית שלה
גם כן רציפה.
\end{itemize}
אם יש הומיאומורפיזם \L{$f:X\to Y$} אומרים ש-\L{$X,Y$} \textbf{הומיאומורפיים}.
שימו לב שהמשמעות של זה היא לא ש-\L{$X,Y$} הם אותו דבר בדיוק, אלא
שהם אותו דבר בדיוק בכל הנוגע \textbf{לתכונות הטופולוגיות} שלהם. הסיפור
הזה עם ספל הקפה שהוא אותו דבר כמו דונאט? פורמלית הוא אומר שהמרחבים
שמתארים את ספל הקפה ואת הדונאט הם הומיאומורפיים. זה מוביל באופן טבעי
לרצון למצוא קריטריונים נוחים לבדיקה מתי שני מרחבים הם הומיאומורפיים
ומתי לא; זה למשל מה ש\textbf{טופולוגיה אלגברית} מתעסקת בו כשהיא מגדירה
בצורה חכמה מאוד מבנים אלגבריים שמתארים אספקטים מסויימים של המרחב הטופולוגי
שנשמרים תחת הומיאומורפיזם.

אני לא הולך לדבר על זה בכלל בשלב הזה.

במקום זה, בואו נראה עוד אפיון של פונקציות רציפות. הגדרתי אותן בעזרת
קבוצות \textbf{פתוחות}, אבל במרחב טופולוגי לקבוצות הסגורות יש תפקיד
חשוב באותה מידה בערך )הרי המשלימה של כל קבוצה פתוחה היא קבוצה סגורה(.
האם אפשר להגדיר רציפות גם בעזרת קבוצות סגורות? ובכן, בקלות:
\begin{itemize}
\item \L{$f:X\to Y$} היא רציפה אם ורק אם לכל קבוצה סגורה \L{$V\subseteq Y$},
המקור \L{$f^{-1}\left(V\right)$} הוא קבוצה סגורה
\end{itemize}
כלומר, בסך הכל לקחתי את ההגדרה הרגילה והחלפתי \textquotedblleft פתוחה\textquotedblright{}
ב\textquotedblright סגורה\textquotedblright . למה שזה יעבוד? תכונות
בסיסיות של משלים שקצת מזכירות את מה שעשיתי בתחילת הפוסט. נניח ש-\L{$V$}
סגורה, אז \L{$V^{c}$} פתוחה. עכשיו, תכף אראה ש-\L{$f^{-1}\left(V^{c}\right)=\left(f^{-1}\left(V\right)\right)^{c}$}
ולכן אם \L{$V$} סגורה, אז \L{$V^{c}$} פתוחה, והרציפות של \L{$f$}
מלמדת ש-\L{$f^{-1}\left(V\right)$} פתוחה, ולכן \L{$\left(f^{-1}\left(V\right)\right)^{c}$}
סגורה. ואת אותו טריק אפשר לעשות גם בכיוון השני של הוכחת השקילות של
ההגדרות רק כשהפעם מתחילים עם קבוצה פתוחה ולא סגורה.

למה \L{$f^{-1}\left(V^{c}\right)=\left(f^{-1}\left(V\right)\right)^{c}$}?
זה כבר סתם להטוטים של תורת הקבוצות: \L{$x\in f^{-1}\left(V^{c}\right)$}
אם ורק אם \L{$f\left(x\right)\in V^{c}$} כלומר אם ורק אם \L{$f\left(x\right)\notin V$}
כלומר אם ורק אם \L{$x\notin f^{-1}\left(V\right)$} כלומר אם ורק אם
\L{$x\in\left(f^{-1}\left(V\right)\right)^{-1}$}. שום דבר חריג פה.

אפשר לסיים פה, אבל למה לא להראות שקילות עוד יותר נחמדה אפילו?
\begin{itemize}
\item \L{$f$} היא רציפה אם ורק אם לכל \L{$A\subseteq X$} מתקיים \L{$f\left(\overline{A}\right)\subseteq\overline{f\left(A\right)}$}
\end{itemize}
בואו נתחיל מלהניח ש-\L{$f$} רציפה ונוכיח ש-\L{$f\left(\overline{A}\right)\subseteq\overline{f\left(A\right)}$}.
כלומר, אני לוקח \L{$y\in f\left(\overline{A}\right)$} ורוצה להראות
ש-\L{$y\in\overline{f\left(A\right)}$}, כלומר ש-\L{$y$} היא נקודת
הצטברות של \L{$f\left(A\right)$}. ראשית, \L{$y\in\overline{f\left(A\right)}$}
פירושו שקיים \L{$x\in\overline{A}$} כך ש-\L{$f\left(x\right)=y$}.
שנית, נקודות הצטברות מוגדרות בעזרת סביבות אז בואו ניקח סביבה \L{$V$}
של \L{$y$} ונוכיח שיש בה איבר אחד לפחות מ-\L{$f\left(A\right)$}.
מכיוון ש-\L{$V$} קבוצה פתוחה ו-\L{$f$} רציפה, \L{$f^{-1}\left(V\right)$}
היא קבוצה פתוחה שבנוסף לכך גם מקיימת \L{$x\in f^{-1}\left(V\right)$},
כלומר היא סביבה של \L{$x$}. מכיוון ש-\L{$x\in\overline{A}$} כל סביבה
של \L{$x$} מכילה נקודה \L{$a\in A$}, כלומר \L{$a\in f^{-1}\left(V\right)$},
כלומר \L{$f\left(a\right)\in V$} והנה מצאנו איבר של \L{$f\left(A\right)$}
ששייך לסביבה \L{$V$}. זה מסיים את הכיוון הזה של ההוכחה.

הכיוון השני מעניין יותר )אותי, בגלל שלי אישית מרגיש שהתכונה של \L{$f\left(\overline{A}\right)\subseteq\overline{f\left(A\right)}$}
\textquotedblleft חלשה יותר\textquotedblright{} מאשר רציפות(. כדי להוכיח
ש-\L{$f$} רציפה, יותר מתבקש להראות שאם \L{$V$} \textbf{סגורה} אז
\L{$f^{-1}\left(V\right)$} סגורה, כי התכונה שכבר יש לי ביד מתעסקת
עם סגור ולכן סביר שיהיה יותר קל לדבר על קבוצות סגורות.

התכונה אומרת ש-\L{$f\left(\overline{A}\right)\subseteq\overline{f\left(A\right)}$},
ואנחנו רוצים לקשר את זה איכשהו אל \L{$V$} ואל \L{$f^{-1}$}. בואו
נסמן \L{$A=f^{-1}\left(V\right)$}, אז \L{$f\left(A\right)=f\left(f^{-1}\left(V\right)\right)\subseteq V$}.
שימו לב שזה \textbf{לא} חייב להיות שוויון, כי בהחלט ייתכן שיש ב-\L{$V$}
איברים שבכלל לא מתקבלים על ידי \L{$f$} ולכן לא יהיו שייכים ל-\L{$f^{-1}\left(V\right)$}
)רק אם \L{$f$} \textbf{על} זה לא יקרה(.

עכשיו, המטרה שלנו היא להוכיח ש-\L{$A$} סגורה, כלומר \L{$A=\overline{A}$}.
בואו ניקח \L{$a\in\overline{A}$} ונוכיח ש-\L{$a\in A$}. עכשיו צריך
איכשהו להשתמש בנתון \L{$f\left(\overline{A}\right)\subseteq\overline{f\left(A\right)}$},
אז בואו נפעיל את \L{$f$} על \L{$a$} ונקבל:

\L{$f\left(a\right)\in f\left(\overline{A}\right)\subseteq\overline{f\left(A\right)}\subseteq\overline{V}=V$}

כאן השתמשתי בכך שראינו כבר \L{$f\left(A\right)\subseteq V$} ולכן
\L{$\overline{f\left(A\right)}\subseteq\overline{V}$}, כלומר שלקחת
סגור לקבוצה זו פעולה \textbf{מונוטונית}. לא הוכחתי את זה קודם אבל
זה די ברור: אם \L{$A\subseteq B$} הן שתי קבוצות כלשהן, אז \L{$\overline{B}$}
היא קבוצה סגורה שמכילה את \L{$A$}, ולכן היא אחת מהקבוצות שמשתתפות
בחיתוך שנותן את \L{$\overline{A}$} כך ש-\L{$\overline{A}\subseteq\overline{B}$}.

אוקיי, אז מה קיבלנו? \L{$f\left(a\right)\in V$} מילולית אומר ש-\L{$a\in f^{-1}\left(V\right)=A$},
וזה בדיוק מה שרצינו. זה מסיים את ההוכחה של הטענה הזו.

לסיום, בואו נראה שתי תוצאות שימושיות על איך להראות שפונקציות הן רציפות
שאצטרך כבר בפוסט הבא אבל עדיף להראות אותן עכשיו כדי לא לקטוע את הזרימה
בו ולהיעזר בכך שכרגע אנחנו עדיין זוכרים את ההגדרה.

ראשית, אם רוצים להוכיח שפונקציה \L{$f:X\to Y$} היא רציפה ישירות מתוך
ההגדרה, לכל תת-קבוצה \L{$V\subseteq Y$} פתוחה צריך להראות ש-\L{$f^{-1}\left(V\right)$}
פתוחה. זה מעיק, כי קבוצות פתוחות יכולות להיות די מסובכות - לכן טוב
לדעת שאפשר להצמצם רק לבסיסים/תת-בסיסים של הטופולוגיה. כלומר: אם \L{$\mathcal{B}$}
הוא תת-בסיס או בסיס של \L{$\left(Y,\mathcal{T}_{Y}\right)$} אז מספיק
להראות ש-\L{$f^{-1}\left(B\right)$} פתוחה ב-\L{$X$} לכל \L{$B\in\mathcal{B}$}
כדי להוכיח ש-\L{$f$} רציפה.

ההוכחה של זה תעבור דרך תכונות תורת-קבוצתיות של פונקציות, כמו שהשתמשתי
קודם ב-\L{$f^{-1}\left(V^{c}\right)=\left(f^{-1}\left(V\right)\right)^{c}$}.
איך \L{$f^{-1}$} משחק עם חיתוכים ואיחודים? ובכן, בדיוק באותה צורה,
באופן האופטימלי ביותר שאפשר:
\begin{itemize}
\item \L{$f^{-1}\left(\bigcup B_{i}\right)=\bigcup f^{-1}\left(B_{i}\right)$}
\item \L{$f^{-1}\left(\bigcap B_{i}\right)=\bigcap f^{-1}\left(B_{i}\right)$}
\end{itemize}
ההוכחה היא שרשרת אם-ורק-אם פשוטה, ממש ברמת תורת הקבוצות של סמסטר ראשון:

\L{$x\in f^{-1}\left(\bigcup B_{i}\right)\iff f\left(x\right)\in\bigcup B_{i}\iff\exists i:f\left(x\right)\in B_{i}\iff$}

\L{$\iff\exists i:x\in f^{-1}\left(B_{i}\right)\iff x\in\bigcup f^{-1}\left(B_{i}\right)$}

ובאופן דומה גם עבור חיתוכים.

המסקנה? אם \L{$\mathcal{B}$} בסיס/תת-בסיס של \L{$Y$} ו-\L{$V\subseteq Y$}
פתוחה, אז מכיוון שאפשר לכתוב את \L{$V$} בעזרת איחודים וחיתוכים של
אברי \L{$\mathcal{B}$} נקבל ש-\L{$f^{-1}\left(V\right)$} הוא איחודים
וחיתוכים של הפעלת \L{$f^{-1}$} על אברי \L{$\mathcal{B}$}, כך שאם
כולם פתוחים נקבל שגם \L{$f^{-1}\left(V\right)$} פתוחה.

שנית, אם \L{$f:X\to Y$} ו-\L{$g:Y\to Z$} הן פונקציות רציפות אז גם
ההרכבה \L{$gf$} היא רציפה )יש כאלו שכותבים הרכבה בתור \L{$f\circ g$}
או \L{$g\circ f$}; אני אוותר על זה כאן. המשמעות של \L{$gf$} היא
הפונקציה שבה קודם מפעילים את \L{$f$} ואז מפעילים את \L{$g$} על התוצאה(.
זו תוצאה סטנדרטית בחשבון אינפיניטסימלי אבל שם מוכיחים את זה עם אפסילון-דלתא;
בואו נעשה את זה עם ההגדרה הכללית יותר. צריך להיות טיפה זהירים עם מה
זה בדיוק אומר \L{$\left(gf\right)^{-1}$} אבל זו תהיה הוכחה די פשוטה.
נתחיל עם \L{$W\subseteq Z$} פתוחה. 

עכשיו, בואו נסמן \L{$V=g^{-1}\left(W\right)$} ו-\L{$U=f^{-1}\left(V\right)$}.
אם \L{$W$} פתוחה, אז \L{$V$} פתוחה כי \L{$g$} רציפה, ולכן \L{$U$}
פתוחה כי \L{$f$} רציפה; אם אני אראה ש-\L{$\left(gf\right)^{-1}=U$},
סיימתי. עכשיו, על פי ההגדרה:

\L{$\left(gf\right)^{-1}\left(W\right)=\left\{ x\in X\ |\ g\left(f\left(x\right)\right)\in W\right\} $}

מצד אחד, אם \L{$x\in U$} אז \L{$f\left(x\right)\in V$} ולכן \L{$g\left(f\left(x\right)\right)\in W$},
כלומר \L{$x\in\left(gf\right)^{-1}\left(W\right)$}.

מצד שני, אם \L{$x\in\left(gf\right)^{-1}\left(W\right)$} ואני מסמן
\L{$y=f\left(x\right)$} אז אני יודע ש-\L{$g\left(y\right)\in W$}
כלומר \L{$y\in g^{-1}\left(W\right)=V$}. קיבלתי ש-\L{$f\left(x\right)\in V$}
ולכן \L{$x\in f^{-1}\left(V\right)=U$} וזה מסיים את ההוכחה.

האם זה כל מה שיש להגיד על פונקציות רציפות? בוודאי שלא, אבל בשביל להגיד
עוד דברים אני צריך עוד סוגי מרחבים לדבר עליהם, ואל זה נגיע בפוסט הבא.
\end{document}
